\section{Reciprocidad theta--Mellin y conservación de la estructura polar}
\label{rec:espectral}

\subsection{Periodización de la red generada}

La red \(\Lambda_t\) del teorema~\ref{rec:thm:red} posee una lectura analítica que conserva su dependencia del registro. Definimos, para \(x>0\),
\begin{equation}
 \Theta_t(x)=\sum_{v\in\Lambda}
       \exp\!\left(-\pi_{\HMT}x\|G_tv\|^2\right).
 \label{rec:eq:theta}
\end{equation}
En esta realización, la coordenada \(\pi_{\HMT}\), previamente generada, se reconoce con la normalización circular de Fourier. Ese reconocimiento se produce después de la construcción de la red y del flujo. Para abreviar la escritura analítica utilizaremos \(\pi\) con ese significado.

En todo intervalo compacto de \(t\), los valores singulares de \(G_t\) están acotados superior e inferiormente por constantes positivas. La red tiene un número \(O(R^{24})\) de puntos en una bola de radio \(R\). La comparación por coronas con una gaussiana demuestra convergencia absoluta y uniforme de \eqref{rec:eq:theta}, y de sus derivadas, cuando \(x\) permanece en un compacto de \((0,\infty)\). El modo \(v=0\) contribuye exactamente uno.

\begin{theorem}[Reciprocidad de la lectura theta]
\label{rec:thm:theta}
Para todo \(x>0\) y \(t\in\RR\),
\begin{equation}
 \Theta_t(x)=x^{-12}\Theta_{-t}(1/x).
 \label{rec:eq:theta-reciproca}
\end{equation}
\end{theorem}
\begin{proof}
Con la convención \(\widehat f(\xi)=\int f(y)e^{-2\pi i\langle y,\xi\rangle}\,dy\), la gaussiana \(f_x(y)=e^{-\pi x\|y\|^2}\) en dimensión veinticuatro tiene transformada
\(\widehat f_x(\xi)=x^{-12}e^{-\pi\|\xi\|^2/x}\).
Periodizamos sobre \(\Lambda_t\). Integrar la periodización sobre un dominio fundamental muestra que el coeficiente de Fourier indexado por \(\xi\in\Lambda_t^*\) es \(\widehat f_x(\xi)/\operatorname{covol}(\Lambda_t)\). Las dos series gaussianas convergen absolutamente; evaluar la expansión en el origen da
\[
 \sum_{v\in\Lambda_t}f_x(v)
 =\operatorname{covol}(\Lambda_t)^{-1}
   \sum_{\xi\in\Lambda_t^*}\widehat f_x(\xi).
\]
Se sustituyen \(\operatorname{covol}(\Lambda_t)=1\) y \(\Lambda_t^*=\Lambda_{-t}\). La demostración conserva el modo constante en ambos miembros.
\end{proof}

La potencia \(x^{-12}\) procede del rango veinticuatro. La inversión del parámetro es consecuencia de la red dual; no se añade como una simetría independiente de la lectura.

\subsection{Continuación meromorfa y ecuación funcional}

Para \(\Re s>12\), sea
\begin{equation}
 \mathcal E_t(s)=\sum_{v\in\Lambda\setminus\{0\}}\|G_tv\|^{-2s},
 \qquad \mathcal L_t(s)=\pi^{-s}\Gamma(s)\mathcal E_t(s).
 \label{rec:eq:epstein}
\end{equation}
El conteo reticular anterior demuestra convergencia absoluta en ese semiplano. La integral gamma, aplicada término a término, produce
\[
 \mathcal L_t(s)=\int_0^\infty(\Theta_t(x)-1)x^{s-1}\,dx.
\]

\begin{theorem}[Descomposición polar y reciprocidad bilateral]
\label{rec:thm:mellin}
La función \(\mathcal L_t\) se prolonga meromórficamente mediante
\begin{align}
 \mathcal L_t(s)={}&\frac1{s-12}-\frac1s\nonumber\\
 &+\int_1^\infty\left[(\Theta_t(x)-1)x^{s-1}
                 +(\Theta_{-t}(x)-1)x^{11-s}\right]\,dx.
 \label{rec:eq:mellin-partida}
\end{align}
Sus únicos polos son simples, en \(s=0,12\), con residuos \(-1,+1\), respectivamente. Se cumple
\begin{equation}
 \mathcal L_t(s)=\mathcal L_{-t}(12-s),\qquad
 \Xi_t(s):=s(s-12)\mathcal L_t(s)=\Xi_{-t}(12-s),
 \label{rec:eq:funcional}
\end{equation}
y \(\Xi_t\) es entera.
\end{theorem}
\begin{proof}
Separamos la integral en \(x=1\). Para \(0<x<1\), la reciprocidad theta da
\[
 \Theta_t(x)-1=x^{-12}-1+x^{-12}(\Theta_{-t}(1/x)-1).
\]
Los dos primeros términos integran a \(1/(s-12)-1/s\). En el tercero, el cambio \(y=1/x\) produce
\(\int_1^\infty(\Theta_{-t}(y)-1)y^{11-s}\,dy\).
La longitud mínima positiva de la red garantiza decaimiento exponencial de \(\Theta_t(x)-1\) en \(x\to\infty\). En cada compacto de \(s\), ese decaimiento domina cualquier potencia de \(x\) y de \(\log x\); por ello las integrales y sus derivadas complejas convergen uniformemente y definen funciones enteras. Los residuos se leen en la parte racional. Sustituir simultáneamente \((t,s)\) por \((-t,12-s)\) intercambia las integrales y deja invariante la parte racional. Multiplicar por \(s(s-12)\) elimina ambos polos.
\end{proof}

El centro \(s=6\) resulta del rango y del exponente \(-2s\) adoptado en \eqref{rec:eq:epstein}. La colección es la completación de Epstein de las redes deformadas. Su variable espectral y su normalización permanecen distintas de las de la función zeta de Riemann; la relación probada entre ambas construcciones es la reciprocidad de sus realizaciones circulares y reticulares.

\begin{corollary}[Conservación polar y componente antisimétrica]
\label{rec:cor:polar}
Las funciones
\[
 \mathcal L_t(s)-\mathcal L_0(s),\qquad
 \mathcal A_t(s)=\frac{\mathcal L_t(s)-\mathcal L_{-t}(s)}2
\]
son enteras. Además,
\[
 \mathcal A_t(12-s)=-\mathcal A_t(s),\qquad \mathcal A_t(6)=0.
\]
\end{corollary}
\begin{proof}
La parte racional de \eqref{rec:eq:mellin-partida} es independiente de \(t\), por lo que se cancela en ambas diferencias. La ecuación funcional cambia el signo de la segunda. Evaluar en su punto fijo da \(2\mathcal A_t(6)=0\).
\end{proof}

La deformación reside, por tanto, en una diferencia entera y conserva exactamente los datos polares. El cero central pertenece a la diferencia bilateral, sin imponer un cero a cada \(\mathcal L_t\). Los polos descritos son singularidades de una transformada de Mellin; la conexión física con el operador de masas se establece más adelante mediante la sección de acción, con su variable y su lector propios.

\subsection{Reconstrucción estable desde el archivo bilateral}
\label{rec:archivo}

La ley de recuperación anterior proporciona, para el desplazamiento cíclico unitario \(S\),
\[
 T_j=\frac{8I+S^j}{9},\qquad D_j=\frac{\sqrt8(S^j-I)}9,
 \quad j=3,4,
\]
\begin{equation}
 m=(D_3K,D_4T_3K)=MK,\qquad LM=P_{11},\qquad\|L\|=\frac94.
 \label{rec:eq:archivo}
\end{equation}
En consecuencia,
\begin{equation}
 u=P_3Lm.
 \label{rec:eq:recuperacion-u}
\end{equation}
Fijadas la coordenatización de incidencia, la red marcada y la razón regional \(\rho\), esta igualdad recupera consecutivamente \(u/\|u\|\), \(g_t\), \(G_t\), \(\Lambda_t\), \(\Theta_t\) y \(\mathcal E_t\). La reconstrucción de esta colección no requiere el terminal \(T_4T_3K\) ni la carga uniforme del registro. Los antecedentes geométricos permanecen incluidos entre los datos fijados.

\begin{theorem}[Estabilidad geométrica y espectral de la recuperación]
\label{rec:thm:estabilidad}
Supongamos \(m'=m+e\), \(\|e\|\le\varepsilon\) y \(9\varepsilon/4<\|u\|\). Sea \(u'=P_3Lm'\), y reconstruyamos los flujos con \(u'/\|u'\|\), sin variar \(\rho,P_3,L\) ni las coordenatizaciones. Si \(a=|t\log\rho|\), entonces
\begin{equation}
 \|g'_t-g_t\|,\ \|G'_t-G_t\|
 \le\frac92\,a e^a\frac{\varepsilon}{\|u\|}.
 \label{rec:eq:estabilidad}
\end{equation}
Con \(H_K=\operatorname{diag}(h_i)\), \(b=|t|\|H'_K-H_K\|\), se tiene
\begin{align}
 b&\le\frac{9a\varepsilon}{2\|u\|},\nonumber\\
 \Theta_t(e^{2b}x)&\le\Theta'_t(x)\le\Theta_t(e^{-2b}x),\label{rec:eq:theta-error}\\
 e^{-2\sigma b}\mathcal E_t(\sigma)&\le\mathcal E'_t(\sigma)
             \le e^{2\sigma b}\mathcal E_t(\sigma),\qquad\sigma>12.
 \label{rec:eq:epstein-error}
\end{align}
\end{theorem}
\begin{proof}
La norma del recuperador da \(\|u'-u\|\le9\varepsilon/4\), que garantiza \(u'\ne0\). La desigualdad triangular y la desigualdad inversa para normas implican
\[
 \left\|\frac{u'}{\|u'\|}-\frac{u}{\|u\|}\right\|
 \le\frac{2\|u'-u\|}{\|u\|}.
\]
Cada exponente está entre \(-a\) y \(a\). Aplicar el teorema del valor medio a la exponencial y tomar la norma diagonal prueba \eqref{rec:eq:estabilidad}; repetir cada entrada dos veces no cambia esa norma. La misma estimación antes de exponenciar da la cota de \(b\).

Como los dos generadores son diagonales en la misma coordenatización,
\[
 e^{-2b}\|G_tv\|^2\le\|G'_tv\|^2\le e^{2b}\|G_tv\|^2.
\]
Aplicar la exponencial decreciente y sumar prueba \eqref{rec:eq:theta-error}. Elevar a \(-\sigma\) y sumar en el dominio de convergencia prueba \eqref{rec:eq:epstein-error}. Las cotas se aplican a todos los vectores, no a una truncación finita de las series.
\end{proof}

El error \(\varepsilon\) se mide en las memorias normalizadas de \eqref{rec:eq:archivo}. Si se conservan los registros racionales \(z=m/\sqrt8\), su error se multiplica por \(\sqrt8\) antes de aplicar la estimación. Los errores independientes en la razón regional o en las coordenatizaciones requieren sus propias cotas. El resultado recupera una colección de geometrías y lecturas analíticas, con un control cuantitativo del archivo que la determina.
